\documentclass[11pt]{article}
\usepackage[margin=1in]{geometry}
\usepackage{amsmath,amssymb,amsthm,bm,hyperref,booktabs,enumitem}
\theoremstyle{plain}
\newtheorem{theorem}{Theorem}\newtheorem{proposition}{Proposition}\newtheorem{corollary}{Corollary}
\theoremstyle{definition}\newtheorem{definition}{Definition}\newtheorem{assumption}{Assumption}
\newcommand{\R}{\mathbb{R}}\newcommand{\E}{\mathbb{E}}\newcommand{\N}{\mathcal{N}}
\newcommand{\Pf}{\mathbb{P}}\newcommand{\I}{\mathbb{1}}\newcommand{\HH}{\mathbb{H}}
\newcommand{\Lpath}{L^\star_{\mathrm{path}}}\newcommand{\Lstr}{L^\star_{\mathrm{struct}}}

\title{\textbf{Structural De-aliasing: Recovering Latent States that a\\ Measurement Channel Cannot Separate}\\[4pt]
\large Minimax rates, measurement design, and a singular local regime}
\author{Robert Richardson}
\date{\today}

\begin{document}\maketitle

\begin{abstract}
A cheap measurement channel---an automated grader, a weak annotator, a machine verifier---reports a noisy label of
a latent scientific state at each node of a known dependency structure. We study the regime in which the channel
\emph{nearly aliases} two scientifically distinct states: their node-level emission distributions coincide or
nearly so, so no amount of the channel, or better calibration of it, separates them. The states nonetheless
participate \emph{differently} in the dependency structure, and the question we answer is quantitative: at what
rate can the structural observations de-alias states the channel cannot? For a node with $d$ conditionally
independent structural views, the effective information combines measurement and structural separation as
$\delta_E^2+d\,\delta_G^2$, with error decaying at the Chernoff exponent between the structural signatures. We lift
this to the field problem of recovering $\{h(Z_v)\}$ on a tree: the belief-propagation marginal-mode estimator is
Bayes, the forest minimax risk has \emph{exactly} the single-gadget Chernoff exponent, and a data-processing
\emph{branch-information sandwich}---an entire branch reveals a node at most as well as its first hidden
neighbour---pins the known-channel field rate on any tree under a mild observability condition, with no spectral
gap or packing argument. Two distinct information axes emerge---local structural richness $d\,\delta_G^2$
(classification) and population size via $\lambda=\sqrt n\,\delta_G^2$ (learning the channel)---yielding a
two-dimensional phase diagram with two separate failure modes. We give a measurement-design law that prices
structural views against replication and a second instrument, and we show a scientific functional $R=h(Z)$ can be
point-identified while the latent state and the confusion matrix are not. Finally, the local weak-identification
regime is genuinely \emph{singular}: an exchangeability symmetry makes the likelihood even in the separation, so
the identifiable \emph{magnitude} $\psi=\delta_G^2$ is regular (estimable at rate $n^{-1/4}$ in $\delta_G$) with a
data-driven boundary-truncated-Gaussian posterior, while the \emph{orientation} of the states is not identified by
the data and is fixed by a single anchor. A natural-language--to--evidence
retrieval problem, where an LLM grader is blind to bridging evidence, is the running illustration and a predicted
corner ($\mathcal I_G\!\to\!0$) of the same theory.

\smallskip\noindent\emph{Keywords:} aliasing; differential misclassification; hidden Markov random fields;
minimax rates; weak identification; singular models; measurement design; partial identification.
\end{abstract}

\section{Introduction}
\paragraph{The problem.} Let a latent state $Z_v\in[K]$ sit at each node $v$ of a known tree or graph, evolving by
a Markov transition $T$ along edges, and let a measurement channel emit $Y_v\sim O_{Z_v}$. Two states $a,b$ are
\emph{node-aliased} when their emission rows coincide, $O_a=O_b$ (exactly or nearly). The scientific target is a
functional $R_v=h(Z_v)$---a decision, not the full state. When $h(a)\neq h(b)$ for an aliased pair, the channel is
blind to a distinction that matters, and no calibration of the channel recovers it. But if $a$ and $b$ occur next
to \emph{different} states, the dependency structure can distinguish them. This paper asks, and answers
quantitatively, \emph{when and how fast}.

\paragraph{A motivating instance.} In multi-hop retrieval, a large-language-model relevance grader is
\emph{bridge-blind}: it grades a passage that supplies a necessary linking entity---but does not itself answer the
query---exactly like an irrelevant passage. Bridge and irrelevant are node-aliased; the bridge, however, sits
adjacent to directly-relevant evidence, while the irrelevant passage does not. Whether corpus structure can
recover relevance the grader misses is exactly a structural de-aliasing problem, and its empirical behavior
(structure helps in the abstract but not once a strong covariate is present) turns out to be a predicted corner of
the theory, not an anomaly.

\paragraph{Contributions.}
\begin{enumerate}[leftmargin=*,itemsep=1pt,topsep=2pt]
\item \textbf{A near-aliasing rate} (\S\ref{sec:rate}). For a node with $d$ conditionally independent structural
views, the two-point Bayes error for the aliased pair decays as $\exp(-\Theta(\delta_E^2+d\,\delta_G^2))$, with the
sharp exponent the shared-$s$ Chernoff information between the structural signatures. This is a \emph{rate} where
the prior literature on aliased-state hidden Markov models gives only identifiability and consistency.
\item \textbf{Field estimation} (\S\ref{sec:field}). Under average Hamming loss the marginal-mode estimator is
Bayes; the forest minimax risk has \emph{exactly} the single-gadget Chernoff exponent; and a data-processing
branch-information sandwich pins the known-channel field rate on \emph{any} tree under a mild observability
condition---no spectral gap or packing. A clean transition separates measurement-dominated, structure-dominated,
and unidentifiable regimes.
\item \textbf{A two-dimensional phase diagram} (\S\ref{sec:phase}). Local structural richness $d\,\delta_G^2$
(can we classify if the channel map were known?) and population information $\lambda=\sqrt n\,\delta_G^2$ (can we
learn the map?) are distinct axes, with two separate failure modes.
\item \textbf{A measurement-design law} (\S\ref{sec:design}): independent channels add Chernoff exponents, so
$n$ replications, $d$ structural views, and $m$ readings of a second instrument contribute
$nC_E+dC_G+mC_2$ (first order), pricing structure against measurement.
\item \textbf{Functional identifiability without parameter identifiability} (\S\ref{sec:func}): the decision
$R=h(Z)$ can be point-identified while the state and the confusion matrix are not.
\item \textbf{A singular local regime} (\S\ref{sec:singular}). By an exchangeability symmetry the local likelihood
is \emph{even} in the separation; the identifiable parameter is $\psi=\delta_G^2$, estimable at rate $n^{-1/4}$ in
$\delta_G$, with a boundary-truncated Gaussian posterior; the orientation of the aliased pair is fixed by an
anchor, not the data.
\end{enumerate}

\paragraph{Relation to existing work.} The qualitative identifiability backbone is known and we use it as
scaffolding: observational equivalence of the structural process is probabilistic bisimulation, equivalently
emission-consistent strong lumpability (Larsen \& Skou, 1991; Kemeny \& Snell, 1960; Buchholz, 1994), and that the
branching neighborhood distinguishes more than any single path is the classical linear-vs-branching-time gap
(Baier et al., 2005). Identifiability of latent tree models from repeated conditionally independent views is
classical (Chang, 1996; Allman, Matias \& Rhodes, 2009), but assumes \emph{distinct} emissions and gives no rates.
Aliased-state HMMs were treated by Weiss \& Nadler (2015) for two states on a line, with identifiability and
consistency but \emph{no rate} and not in the near-aliased regime. Our contribution is the quantitative theory
in the near-aliased regime---rates, minimax field recovery, design, and the singular local asymptotics---together
with the aliasing framing that the distinct-emission literature does not cover. Functional recovery under
non-identification continues the tradition of estimable functions (Rao, 1973) and partial identification (Manski,
2003).

\section{Model and aliasing depths}\label{sec:model}
On a tree, $\{Z_v\}$ is a Markov random field with edge transition $T$ and the measurement channel emits
$Y_v\mid Z_v=c\sim O_c$. Optional cheap covariates $X_v$ may accompany each node. The target is $R_v=h(Z_v)$.
\begin{definition}[Path and structural signatures]
The \emph{path signature} $P_\ell(c)=\mathcal L(Y_w\mid Z_v=c)=(T^\ell O)_c$ is the law of one measurement at
distance $\ell$; the \emph{structural signature} $Q_\ell(c)=\mathcal L(Y_{B_\ell(v)}\mid Z_v=c)$ is the joint law
of the whole radius-$\ell$ neighborhood. The corresponding aliasing depths are
$\Lpath=\min\{\ell:P_\ell(a)\neq P_\ell(b)\}$ and $\Lstr=\min\{\ell:Q_\ell(a)\neq Q_\ell(b)\}$, with
$\Lstr\le\Lpath$.
\end{definition}
Node-level aliasing ($O_a=O_b$) makes $\Lstr,\Lpath>0$; heterophily (aliased states have different neighbors) makes
them finite. The joint signature is strictly richer than the marginal: because siblings sharing a hidden parent are
conditionally dependent, $Q_\ell$ can separate states whose every path marginal $P_\ell$ coincides---a de Finetti /
overdispersion effect that can give $\Lstr<\Lpath$ (Prop.~\ref{prop:branch}). Write $\delta_E^2=H^2(O_a,O_b)$ and
$\delta_G^2=H^2\big((TO)_a,(TO)_b\big)$ for the squared-Hellinger node and structural separations.

\section{The near-aliasing rate}\label{sec:rate}
\begin{theorem}[Near-aliasing rate]\label{thm:rate}
Test $Z_v=a$ vs $Z_v=b$ from the node measurement and $d$ conditionally independent radius-$L$ views. With
per-channel Chernoff divergences $C_E(s),C_G(s)$ and Bhattacharyya coefficients $\rho_E,\rho_G$, the equal-prior
Bayes error obeys $\tfrac14(\rho_E\rho_G^d)^2\le P_e\le\tfrac12\rho_E\rho_G^d$; the exact consistency condition is
$-\log\rho_E-d\log\rho_G\to\infty$, which \emph{near the alias} reads $\delta_E^2+d\,\delta_G^2\to\infty$. The sharp
exponent is the shared-tilt Chernoff information $\max_s\{C_E(s)+d\,C_G(s)\}$, equal to
$\tfrac12(\delta_E^2+d\,\delta_G^2)(1+o(1))$ near the alias; under a nonlattice, nondegenerate-variance condition
the prefactor is $d^{-1/2}$ (for a fixed branch channel), and $(d\,\delta_G^2)^{-1/2}$ if the channel itself
approaches aliasing with $d$.
\end{theorem}
\begin{proposition}[Branching de-aliasing]\label{prop:branch}
There is a finite two-layer model with node-aliased $a,b$ whose every constituent single measurement is identically
distributed yet whose two-sibling joint differs; the de-aliasing exponent is the positive Chernoff information of
the joint signatures, though the per-measurement Chernoff information is zero. (This is the linear-vs-branching gap;
we state only the finite two-layer form, not a homogeneous-chain $\Lpath=\infty$.)
\end{proposition}

\section{Field estimation: oracle risk, exact forest exponent, and the branch bottleneck}\label{sec:field}
We keep the \emph{oracle} problem (channel $\vartheta{=}(T,O)$ known, latent field unknown) and the \emph{minimax}
problem ($\vartheta$ unknown) sharply distinct. Under average Hamming loss the Bayes rule is the vector of marginal
posterior modes, computed exactly by belief propagation, so the oracle risk is the average nodewise two-point
Bayes error $P_e(v)$. Let $\Phi(x)\asymp e^{-cx}$.
\begin{theorem}[Forest risk, exact exponent; $\vartheta$ known]\label{thm:forest}
On the forest of $n$ disjoint depth-1 star gadgets, $P_e\le\mathcal R_{\mathrm{oracle}}\le 2P_e$, where $P_e$ is the
single-gadget two-point Bayes error; hence the forest risk has \emph{exactly} the single-gadget Chernoff exponent
$\max_s\{C_E(s)+dC_G(s)\}$, not merely an $\exp(-\Theta)$ rate.
\end{theorem}
\begin{theorem}[Branch-information sandwich and the known-channel field rate]\label{thm:sandwich}
Writing each branch as the transition row passed through an observation channel, data processing gives, for every
tilt $s$, $C_s\big((TO)_a,(TO)_b\big)\le C_s\big(G_a^{(\infty)},G_b^{(\infty)}\big)\le C_s(T_a,T_b)$: an entire
branch reveals the centre state at most as well as the first hidden neighbour. Under the observability condition
$C_s(T_a,T_b)\le K\,C_s((TO)_a,(TO)_b)$, therefore, the per-node Bayes error on \emph{any} bounded-degree tree
satisfies $-\log P_e(v)=\Theta(\delta_E^2+d\,\delta_G^2)$, and the oracle field risk is
$\mathcal R_{\mathrm{oracle}}\asymp\Phi(\delta_E^2+d\,\delta_G^2)$---with no spectral-gap, packing, or Assouad
argument. If observability fails, deeper observations carry qualitatively new information: exactly the
branching-de-aliasing regime of Prop.~\ref{prop:branch}.
\end{theorem}
\noindent For \emph{unknown} $\vartheta$, the minimax risk is $\mathcal R_{\mathrm{oracle}}$ plus the singular
estimation remainder governed by $\lambda=\sqrt n\,\delta_G^2$ (\S\ref{sec:singular}). Three regimes:
\emph{measurement-dominated} ($\delta_E^2\gg d\delta_G^2$), \emph{structure-dominated} ($\delta_E=0$), and
\emph{unidentifiable} ($\delta_E^2+d\delta_G^2\to0$, risk $\to\tfrac12$).

\section{Two information axes: an $(n,d)$ phase diagram}\label{sec:phase}
Local structural richness $d\,\delta_G^2$ governs whether the neighborhood can classify a node \emph{if the
channel map were known}; population information $\lambda=\sqrt n\,\delta_G^2$ governs whether $n$ nodes suffice to
\emph{learn} that map (\S\ref{sec:singular}). De-aliasing fails in two distinct ways---too little local structure,
or enough structure but too little population information to learn how it maps to the states---and succeeds only
when both are large.

\section{Measurement design}\label{sec:design}
\begin{corollary}[First-order design law]\label{cor:design}
Independent channels add error exponents, so a design with $n$ replicated node measurements, $d$ structural views,
and $m$ readings of a second instrument (per-reading information $C_2$) has effective information equal to the
shared-$s$ exponent $\max_s\{nC_E(s)+dC_G(s)+mC_2(s)\}$, whose near-alias first-order value is
$nC_E+dC_G+mC_2$. Iso-information surfaces give the substitution rates between structure, replication, and a
second instrument; near the alias $C_E\to0$ and structure trades directly against the second instrument.
\end{corollary}

\section{Functional recovery without parameter identification}\label{sec:func}
\begin{theorem}[Functional identifiability]\label{thm:func}
For a fixed oriented parameterization, $R_v=h(Z_v)$ is identified from the radius-$L$ neighborhood iff every pair
$a,b$ with $h(a)\neq h(b)$ is structurally de-aliased; equivalently, the recoverable functionals are those constant
on the observational-equivalence (emission-consistent lumpable) classes, and such an $R$ can be identified even
when $Z$ and $O$ are not. In an \emph{unoriented} model, identification additionally requires invariance under all
label automorphisms preserving the observed law (or an anchor). Distinguishing the classes is model-level
identifiability---not exact recovery of an individual latent draw, whose error is $P_e>0$ (\S\ref{sec:rate}).
\end{theorem}

\section{Unknown parameters and the singular local regime}\label{sec:singular}
When $(T,O)$ are estimated from the same field, the plug-in risk is the oracle floor plus an estimation remainder
$\varrho_n$ governed by $\lambda=\sqrt n\,\delta_G^2$: negligible as $\lambda\to\infty$ (plug-in asymptotically
free), of the same order as the structural gain when $\lambda=O(1)$. The reason is that the parameter separating
the aliased pair is \emph{singular} at the alias.
\begin{proposition}[Singular structure]\label{prop:singular}
Let $\eta$ parametrize the separation ($\eta=0$ full alias, $\delta_G\propto\eta$). By exchangeability of the
aliased pair the likelihood is even, $L(\eta)=L(-\eta)$, so the score at $\eta=0$ vanishes and the information for
$\eta$ is degenerate. The identifiable parameter is $\psi=\eta^2\ (\propto\delta_G^2)$, regular and
$\sqrt n$-estimable; $\eta$ is estimable only at rate $n^{-1/4}$; the invariant local index is
$\lambda=\sqrt n\,\delta_G^2$; and the \emph{sign} of $\eta$---the orientation of the pair, hence $R=h(Z)$---is
not identified by the data.
\end{proposition}
\noindent \emph{Two stages.} Structural data identify the de-aliasing \emph{magnitude} $|\eta|=\sqrt\psi$ and the
\emph{unordered} pair; a single anchor (or the prior) supplies the \emph{orientation} $\mathrm{sign}(\eta)$ and
thus $R=h(Z)$. Orientation costs $O(1)$ external information regardless of $n$: structure identifies how much the
states differ, and very little side information suffices to orient them.
\begin{proposition}[Boundary posterior]\label{prop:boundary}
By a local-asymptotic-normality expansion in $\psi$, along $\psi_n=\lambda_0/\sqrt n$ the posterior for the
localized magnitude $h=\sqrt n\,\psi$ converges to $\N(\Delta/I_\psi,I_\psi^{-1})$ truncated to $[0,\infty)$, with
$\Delta\sim\N(\lambda_0 I_\psi,I_\psi)$. The magnitude posterior is thus \emph{data-driven for every} $\lambda_0$
(at $\lambda_0=0$ it is a data-driven truncated Gaussian, \emph{not} the prior); the permanently prior/anchor-fixed
object is the \emph{orientation} (sign), a Bernoulli. Hence $|\eta|$ lives on the $n^{-1/4}$ scale, while signed
$\eta$ is not estimable at any rate absent orientation.
\end{proposition}

\section{Illustration: bridge-blind retrieval}\label{sec:illus}
Define the incremental structural information $\mathcal I_G=I(R_v;Y_{N(v)}\mid Y_v,X_v)$. Under log loss the
optimal-risk reduction from the neighborhood equals $\mathcal I_G$, and by conditional Pinsker the improvement in
any bounded-loss Bayes risk is at most $\sqrt{\mathcal I_G/2}$. In the retrieval problem, without a strong
covariate $\mathcal I_G$ is large and structure recovers the grader's missed bridges; once a calibrated retriever
covariate is present it is nearly sufficient for relevance, $\mathcal I_G\!\to\!0$, and \emph{no} structural
procedure can materially help---the observed null is the theorem's predicted corner, not a failure of method.

\section{Numerical evidence}\label{sec:num}
Every result is checked by Monte-Carlo simulation (no tuning). Table~\ref{tab:num} collects the confirmations and,
where the theory leaves a constant, the value the simulation pins.
\begin{table}[h]\small
\caption{Simulation checks of each result.}\label{tab:num}
\begin{center}
\begin{tabular}{@{}p{0.31\textwidth}p{0.62\textwidth}@{}}
\toprule
Result & What the simulation shows \\ \midrule
Near-aliasing rate (Thm~\ref{thm:rate}) & shared-tilt exponent $0.845$ at $s^\star{=}0.57$; the separate-max $0.878$ is a strict over-estimate; near-alias boundary $\delta_E^2+d\delta_G^2$ exact \\[2pt]
Branching de-aliasing (Prop~\ref{prop:branch}) & per-measurement Chernoff $0$, joint rate $\gamma{=}0.089>0$ ($\Lstr{=}2$): dependence alone de-aliases \\[2pt]
Field rate (Thms~\ref{thm:forest},\ref{thm:sandwich}) & measurement- and structure-dominated regimes de-alias, aliased-$T$ floors at chance; on a disjoint-neighborhood antichain the field rate equals the local Chernoff rate \\[2pt]
Deep-subtree information & $C_G^{(\ell)}$ grows with depth and saturates to $C_G^{(\infty)}\approx0.122$ (star $C_G^{(1)}\approx0.116$); the forest/gadget class makes the star exponent \emph{exactly} minimax \\[2pt]
Connected-tree lower bound & Assouad on a distance-$2L$ antichain matches the oracle upper bound at the \emph{rate} level (exact constant open) \\[2pt]
Singular local regime (Prop~\ref{prop:singular}) & $\lambda=\sqrt n\,\delta_G^2$ is the invariant index; $\Pf(\mathrm{sign}\,\eta{=}{+})$ moves $0.50\!\to\!0.97$ as anchors go $0\!\to\!10$ \\[2pt]
Boundary posterior (Prop~\ref{prop:boundary}) & an $n$-stable truncated Gaussian set by $\lambda$: half-normal at $\lambda{=}0$, Bernstein--von Mises as $\lambda\to\infty$ \\[2pt]
Plug-in remainder (\S\ref{sec:singular}) & excess risk vanishes as $\lambda\to\infty$ and stays flat while $\sqrt n\,\delta_G$ doubles, confirming $\lambda=\sqrt n\,\delta_G^2$ ($\delta_G\sim n^{-1/4}$) \\
\bottomrule
\end{tabular}
\end{center}
\end{table}

Two points deserve emphasis. The deep-subtree simulation shows per-node structural information grows with depth and
saturates to a finite $C_G^{(\infty)}$ just above the radius-1 value, so the star rate is a conservative bound on a
general tree, while on the forest/gadget class (radius-1, independent neighborhoods) the star Chernoff exponent is
\emph{exactly} minimax. And the unknown-parameter simulation refines the invariant weak-identification index to
$\lambda=\sqrt n\,\delta_G^2$, unifying the field-minimax estimation remainder with the singular local-Bayes regime.

\paragraph{A limitation the simulations make concrete.} Estimating the channel $(T,O)$ from \emph{short} dependency
chains is hard: a chain-EM estimator separates the aliased emission cleanly only as chain length grows, so in the
retrieval instance (chains of $2$--$4$ hops) finite-structure channel identification is weak, even though the
\emph{oracle-channel} de-aliasing rate is unaffected. The theory bites hardest, then, when the channel is known or
estimable; the residual hard problem is joint channel-and-state learning from short structures, governed by the
singular local regime (\S\ref{sec:singular}).

\section{Discussion}
Aliasing turns latent-state recovery into a genuinely two-dimensional problem---local structure versus population
information---with a singular local regime at its boundary. The open items are constants rather than existence:
the exact within-$\Theta$ field constant, the exact constant in the connected-tree lower bound (the rate is
already matched, \S\ref{sec:num}), and the general-tree extension of the boundary-posterior scale
$I_\psi(0)^{-1/2}$ (its shape is proved for the independent-star experiment). The framework is not specific to
language models; it applies wherever an automated or weak measurement channel aliases scientifically distinct
states that behave differently in a known dependency structure.

\appendix
\section{Proof sketches}
\emph{Theorem~\ref{thm:rate}.} Conditional independence of the node emission and the $d$ branches given $Z_v$ makes
the joint law a product; Hellinger affinities multiply and Chernoff exponents add under a shared tilt; Bhattacharyya
and Le Cam give the two-sided bound and Chernoff's theorem the sharp exponent.\quad
\emph{Theorem~\ref{thm:forest}.} Uniform prior on the centre states gives Bayes risk $P_e$ (lower); independent
per-gadget likelihood-ratio tests give $\le 2P_e$ (upper); hence the exact single-gadget Chernoff exponent.\quad
\emph{Theorem~\ref{thm:sandwich}.} Each branch is a transition row through an observation channel, so the branch
Chernoff information is sandwiched by data processing between the immediate-view and the transition-row values;
observability makes both fences $\Theta(\delta_E^2+d\delta_G^2)$.\quad
\emph{Proposition~\ref{prop:singular}.} The label-swap symmetry gives the even likelihood; reparametrizing to
$\psi=\eta^2$ restores regularity, whence the $n^{-1/4}$ rate; the sign is a global relabeling the data cannot
see.\quad
\emph{Proposition~\ref{prop:boundary}.} LAN in $\psi$ with score $S=a/p_0$ and information $I_\psi$; Le~Cam's third
lemma shifts the local score; a locally flat prior yields the truncated-Gaussian posterior for $h=\sqrt n\psi\ge0$.

\smallskip\noindent\footnotesize\emph{Note.} Full proofs are collected in the companion \texttt{paperB\_JASA\_proofs.tex}; a working document
records the empirical provenance and the extended derivations.
\normalsize
\end{document}
